\chapter{テンソル}
    \section{テンソル記法}
        \subsection{例}
            $A,B$を行列，$\vx,\vy$をベクトルとする。
            \begin{enumerate}
                \item $(A\vx)_i = A_{i,j}x_j$
                \item $(AB)_{i,j} = A_{i,k}B_{k,j}$
                \item $(AB\vx)_i = (AB)_{i,k}x_k = A_{i,j}B_{j,k}x_k$
                \item $(\vx\times\vy)_i = \varepsilon_{i,j,k}x_j y_k$
            \end{enumerate}
        \subsection{Levi-Civita記号の性質：\texorpdfstring{$\sum_{a,i,j,k,l}\varepsilon_{a,i,j}\varepsilon_{a,k,l} = \sum_{i,j}(\delta_{i,k}\delta_{j,l} - \delta_{i,l}\delta_{j,k})$}{2 個のεの積と，δの積和}}
            \begin{shadebox}
                \begin{equation}
                    \sum_{a,i,j,k,l}\varepsilon_{a,i,j}\varepsilon_{a,k,l} = \sum_{i,j}(\delta_{i,k}\delta_{j,l} - \delta_{i,l}\delta_{j,k})
                \end{equation}
            \end{shadebox}
            導出は直接比較で得られる。ここでは考え方を示す。
            \par
            左辺のうち非零となる項に於いては，$i,j$の組に対して$a$が唯一に決まる。
            そして1つの$a$に対して$i,j$の組が2つ存在する。
            +1で寄与する項は$(k,l)=(i,j)$であり，-1で寄与する項は$(k,l)=(j,i)$である。
            この時点で右辺が思い付く。
            これが$i=j$のときに0となればよい。
            実際にそうであることは容易に確認できる。